\section{Effective coefficients and certified constants}\label{tks:sec:constants}
At any fixed weight $S$, the product \eqref{tks:eq:P} has a finite polynomial approximation. Enumerate partitions with parts at least $2$ and size at most $S$. Truncate each $d_j$ at weight $S$, replace its reciprocal by
\[
(1-d_j)^{-1}=\sum_{k=0}^{\lfloor S/2\rfloor}d_j^k
\quad\text{modulo terms of weight greater than }S,
\]
and multiply for $2\leq j\leq J$. Every operation uses rational arithmetic. Denote the resulting coefficients by $q_{\rho,J}$ and the root-weighted sums by $H_{s,J}$. Also put
\[
C_J=\prod_{j=2}^{J}(1-1/j!)^{-1},\qquad
R_s^*=\max_{\substack{\rho\vdash s\\\min\rho\geq2}}r(\rho),
\]
with $R_0^*=1$ and an empty maximum interpreted as $0$.

\begin{proposition}\label{tks:prop:tails}
For $J\geq\max(2,s)$,
\begin{align}
0\leq H_s-H_{s,J}&\leq84R_s^*4^{s-J},\label{tks:eq:geomconstanttail}\\
0\leq C-C_J&\leq\frac{6(J+2)}{(J+1)(J+1)!}.\label{tks:eq:Ctail}
\end{align}
For $s\geq2$ and $J\geq\max(s,4)$, the stronger fixed-weight estimate is
\begin{equation}\label{tks:eq:factorialconstanttail}
0\leq H_s-H_{s,J}
 \leq24sR_s^*4^s\frac{(J+1)^s}{(J+1)!}.
\end{equation}
Consequently each $H_s$ is computable to any prescribed absolute precision by a finite algorithm with a proved stopping rule.
\end{proposition}
\begin{proof}
At $t_0=1/4$, expansion of \eqref{tks:eq:Hjt} gives
$\mathcal H_j(1,t_0)\leq e^3 4^{-j}$. Set
\[
\delta_j=d_j(t_0^2,t_0^3,\ldots)
 =\frac{\mathcal H_j(1,t_0)-1/j!}{1-1/j!}\leq2e^3 4^{-j}.
\]
Thus $\sum_{j>J}\delta_j<14\cdot4^{-J}$ and
$P(t_0)\leq G(1,t_0)\leq C_0<6$. Positivity gives
\[
P(t_0)-P_J(t_0)
 =P(t_0)\left(1-\prod_{j>J}(1-\delta_j)\right)
 \leq P(t_0)\sum_{j>J}\delta_j.
\]
Bounding a weight-$s$ coefficient by $4^s$ times this positive total and then multiplying by the maximal root count proves \eqref{tks:eq:geomconstanttail}. The same argument for $C$, using $C\leq6$ and
\[
\sum_{j>J}\frac1{j!}\leq\frac{J+2}{(J+1)(J+1)!},
\]
proves \eqref{tks:eq:Ctail}.

For the refinement, truncate each $d_j$ at weight $s$ before evaluating at $t_0$. This preserves the desired coefficient and can only decrease the positive product. At each weight $k$, the sum of $1/z_\rho$ over types without $1$-parts is at most $1$: it is the proportion of fixed-point-free permutations of $k$. For $j\geq s$, the truncated factor input is therefore at most $2s j^s/j!$. If $j\geq J+1$ and $J\geq\max(s,4)$, then
\[
\frac{(j+1)^s/(j+1)!}{j^s/j!}
 =\frac{(1+1/j)^s}{j+1}\leq\frac e{J+2}<\frac12.
\]
The tail is at most $4s(J+1)^s/(J+1)!$. Multiplication by the product bound $6$, the coefficient factor $4^s$, and $R_s^*$ proves \eqref{tks:eq:factorialconstanttail}.
\end{proof}

\subsection{The first nonzero support coefficients}
Define
\begin{equation}\label{tks:eq:Dk}
D_k=\sum_{j\geq\max(2,k)}\frac{\fall jk}{j!-1},\qquad
v_{j,2}=\frac{\fall j2}{2(j!-1)},\qquad
v_{j,3}=\frac{\fall j3}{3(j!-1)}.
\end{equation}
Direct extraction from \eqref{tks:eq:P} gives
\begin{align}
H_3&=D_3/3,\label{tks:eq:H3}\\
H_4&=D_4/4+(D_2/2)^2+\sum_{j\geq2}v_{j,2}^2,\label{tks:eq:H4}\\
H_5&=D_5/5,\label{tks:eq:H5}\\
H_6&=2D_6/9+2(D_3/3)^2+2\sum_{j\geq3}v_{j,3}^2.\label{tks:eq:H6}
\end{align}
At support $3$ only $(3)$ contributes, with root count $1$; at support $4$ only $(2,2)$ contributes, with root count $2$; at support $5$ only $(5)$ contributes; and at support $6$ only $(3,3)$ contributes, with root count $4$. For example, the coefficient of $u_2^2$ receives both a term from $d_j$ and a term from $d_j^2$ within a single factor, as well as terms from different factors. These repeated-factor contributions are exactly the extra squared sums in \eqref{tks:eq:H4} and \eqref{tks:eq:H6}.

\subsection{Forty decimal interval certificates}
The following are enclosing intervals, not merely rounded floating-point approximations. Every displayed interval has width $10^{-40}$. The finite sums and products use $J=70$, exact rational arithmetic, and the proved tails \eqref{tks:eq:Ctail} and \eqref{tks:eq:factorialconstanttail}; outward decimal rounding is by integer division. The package reproduces the calculation.
\begin{center}
\small
\begin{tabular}{cl}
\toprule
Constant & Certified lower and upper endpoints\\
\midrule
$C$ & $[\,2.5294774720791526481801161542539542411787,$\\
    & $\phantom{[\,}2.5294774720791526481801161542539542411788\,]$\\[3pt]
$H_3$ & $[\,0.9887340176119067106765900015027659522294,$\\
    & $\phantom{[\,}0.9887340176119067106765900015027659522295\,]$\\[3pt]
$H_4$ & $[\,6.0121093539050841781399993190129528805222,$\\
    & $\phantom{[\,}6.0121093539050841781399993190129528805223\,]$\\[3pt]
$H_5$ & $[\,0.5456358974558364595261891034594968055729,$\\
    & $\phantom{[\,}0.5456358974558364595261891034594968055730\,]$\\[3pt]
$H_6$ & $[\,3.1846595830270382795159539520700360539731,$\\
    & $\phantom{[\,}3.1846595830270382795159539520700360539732\,]$\\
\bottomrule
\end{tabular}
\end{center}
The constants need not be rational; the finite approximants and the rigorous error bounds are rational. The certification of these constants is logically separate from an effective remainder estimate for $a_n$: no numerical starting index for an asymptotic theorem follows from this table alone.
\begin{writenote}[The constants recomputed, 7 October 2026]
The write extracted $q_\rho$ for all types of weight at most~6 from its own
truncation of~\eqref{tks:eq:P} at $J=60$ (exact fractions) and evaluated the
closed forms~\eqref{tks:eq:H3}--\eqref{tks:eq:H6} at 50 digits; the two agree
to 50 digits, and every displayed interval contains the recomputed value:
\begin{center}
\begin{tabular}{@{}ll@{}}
$C$ & $2.529477472079152648180116154253954241178702\ldots$\\
$H_3$ & $0.988734017611906710676590001502765952229404\ldots$\\
$H_4$ & $6.012109353905084178139999319012952880522219\ldots$\\
$H_5$ & $0.545635897455836459526189103459496805572989\ldots$\\
$H_6$ & $3.184659583027038279515953952070036053973163\ldots$
\end{tabular}
\end{center}
\end{writenote}
